\section*{Capítulo \ref{ch:b2:ellingham} --- Diagramas de Ellingham y metalurgia}

\begin{solution}{exo:b2:ellingham:1}
\ce{4Cu + O2 -> 2Cu2O}; \ce{4/3Al + O2 -> 2/3Al2O3}; \ce{2C + O2 -> 2CO};
\ce{C + O2 -> CO2}.
\end{solution}

\begin{solution}{exo:b2:ellingham:2}
$\Delta_r H^\circ = 2(-350.46) = \qty{-700.92}{kJ}$; $\Delta_r S^\circ = 2(43.65)
- 2(41.63) - 205.15 = \qty{-201.1}{J/K}$. Así, $\Delta_r G^\circ = -700.9 +
0.2011\,T$ (\unit{kJ}, $T$ en \unit{K}), válida hasta el punto de fusión del cinc.
\end{solution}

\begin{solution}{exo:b2:ellingham:3}
A \qty{1500}{K} la recta de \ce{Cr2O3} está cerca de \qty{-495}{kJ}. Por debajo de ella, y por tanto
capaces de reducir el óxido de cromo: silicio, titanio, aluminio, magnesio y
calcio. Por encima, incapaces: cobre, hierro y cinc. El carbono, a \qty{-488}{kJ}, queda
justo por encima: solo reduce el óxido de cromo a una temperatura algo mayor,
y da entonces un cromo cargado de carburo, una de las razones por las que el cromo para
aleaciones libres de carbono se obtiene por aluminotermia.
\end{solution}

\begin{solution}{exo:b2:ellingham:4}
La pendiente es $-\Delta_r S^\circ$, gobernada por la variación de la cantidad de gas.
\ce{C + O2 -> CO2}: un mol de gas da uno, $\Delta_r S^\circ \approx 0$,
horizontal. \ce{2C + O2 -> 2CO}: uno da dos, $\Delta_r S^\circ > 0$, la
recta desciende. \ce{2CO + O2 -> 2CO2}: tres dan dos, $\Delta_r S^\circ < 0$, la
recta sube, como las de los metales.
\end{solution}

\begin{solution}{exo:b2:ellingham:5}
$\Delta_r G^\circ(\qty{600}{K}) = -181.6 + 600 \times 0.2165 = \qty{-51.7}{kJ}$,
$p_{\ce{O2}} = p^\circ\exp(-51\,710/(8.314 \times 600)) = \qty{3.2e-5}{bar}$. El aire
contiene \qty{0.21}{bar}, mucho más: el mercurio se oxida a \qty{600}{K} (lentamente:
así se preparó por primera vez el óxido), y el óxido solo se descompone
por encima de unos \qty{730}{K}.
\end{solution}

\begin{solution}{exo:b2:ellingham:6}
$\Delta_r G^\circ = -1582.3 - (-743.5) = \qty{-838.8}{kJ}$. Los reactivos son
dos sólidos en contacto solo en la superficie de sus granos, y la reacción
tiene una energía de activación elevada: hay que encenderla localmente (con una mecha de magnesio);
después, el calor que libera la mantiene.
\end{solution}

\begin{solution}{exo:b2:ellingham:7}
La recta de \ce{2C + O2 -> 2CO} corta la de \ce{2Fe + O2 -> 2FeO} cerca de
\qty{1040}{K}; por encima, \ce{FeO + C -> Fe + CO} tiene una \omterm{def:b2:reaction-free-energy:reaction-gibbs}{energía de Gibbs estándar
de reacción} negativa.
\end{solution}

\begin{solution}{exo:b2:ellingham:8}
$\Delta_r G^\circ = 2(-209.1) - (-396.0) = \qty{-22.2}{kJ}$, $K^\circ =
\exp(22\,200/(8.314 \times 1100)) = 11.3$. Entonces $x^2/(1 - x) = 11.3$, $x =
0.92$. Al enfriar a \qty{700}{K} el equilibrio retrocede hacia el \ce{CO2}
($x = 0.016$): el monóxido de carbono se dismuta, \ce{2CO -> C + CO2}, y
se deposita hollín, lentamente salvo que una superficie como el hierro lo catalice.
\end{solution}

\begin{solution}{exo:b2:ellingham:9}
La recta del agua (tres moles de gas dan dos) sube con menos pendiente que la
recta \ce{CO}/\ce{CO2}, también de tres a dos pero con una pérdida de entropía mayor. Las
dos se cortan cerca de \qty{1100}{K}: por debajo, la recta \ce{CO}/\ce{CO2} es más baja y
el monóxido de carbono es el reductor más fuerte; por encima, la recta del agua es más baja y
gana el hidrógeno. Es el mismo enunciado que el \omterm{def:b2:equilibrium-shifts:shift}{desplazamiento del equilibrio} del gas de agua
\ce{CO + H2O <=> CO2 + H2} hacia la izquierda a alta temperatura.
\end{solution}

\begin{solution}{exo:b2:ellingham:10}
Por mol de \ce{O2}, es decir, para \ce{2MgO + Si -> SiO2 + 2Mg(g)}: $\Delta_r
G^\circ = -643.7 - (-845.5) = \qty{+201.8}{kJ}$. Con el silicio y la sílice de
actividad 1, $\Delta_r G = \Delta_r G^\circ + RT\ln(p_{\ce{Mg}}/p^\circ)^2$, negativa
cuando $p_{\ce{Mg}} < p^\circ\exp(-201\,800/(2 \times 8.314 \times 1500)) =
\qty{3e-4}{bar}$. Un horno de vacío bombea el vapor de magnesio a medida que
se forma y lo condensa sobre una superficie fría, de modo que su presión nunca alcanza
ese valor; la cal añadida a la carga fija la sílice y rebaja su actividad,
lo que ayuda todavía más.
\end{solution}

\begin{solution}{exo:b2:ellingham:11}
Cuatro especies, una reacción, tres fases (dos sólidos y el gas): $v = 4 - 1 +
2 - 3 = 2$. A temperatura y presión dadas la razón $p_{\ce{CO}}/p_{\ce{CO2}}$
queda fijada: el gas no puede ceder todo su monóxido de carbono al óxido de hierro,
y el gas de tragante aún contiene mucho, que se quema para precalentar el
aire.
\end{solution}

\begin{solution}{exo:b2:ellingham:12}
$\log K = 2(0.34 + 0.41)/0.059 = 25.4$, $K = \num{3e25}$. With
$[\ce{Fe^2+}] = \qty{1.0}{mol/L}$, $[\ce{Cu^2+}] = 1/K \approx
\qty{4e-26}{mol/L}$: el cobre se elimina por completo.
\end{solution}

\begin{solution}{pb:b2:ellingham:1}
\textbf{1.} $\Delta_r H^\circ = \qty{-700.92}{kJ}$, $\Delta_r S^\circ =
\qty{-201.1}{J/K}$.
\textbf{2.} $\Delta_r G^\circ = -700.9 + 0.2011\,T$ (\unit{kJ}), por debajo de
\qty{692.7}{K}.
\textbf{3.} $\Delta_{\text{fus}}S = 7320/692.7 = \qty{10.6}{J/(K.mol)}$;
$\Delta_{\text{vap}}S = 115\,300/1180.2 = \qty{97.7}{J/(K.mol)}$.
\textbf{4.} Cinc líquido: $\Delta_r H^\circ = -700.92 - 2(7.32) = \qty{-715.6}{kJ}$,
$\Delta_r S^\circ = -201.1 - 2(10.6) = \qty{-222.2}{J/K}$: $\Delta_r G^\circ =
-715.6 + 0.2222\,T$.
\textbf{5.} Cinc gaseoso: $\Delta_r H^\circ = -715.6 - 2(115.3) = \qty{-946.2}{kJ}$,
$\Delta_r S^\circ = -222.2 - 2(97.7) = \qty{-417.7}{J/K}$: $\Delta_r G^\circ =
-946.2 + 0.4177\,T$.
\textbf{6.} Pendientes $0.201$, $0.222$ y $0.418$ \unit{kJ/K}. Por encima del punto de
ebullición la reacción consume tres moles de gas en lugar de uno: la pérdida de
entropía, y con ella la pendiente, se duplica.
\textbf{7.} A \qty{692.7}{K}: $-700.9 + 139.3 = -715.6 + 153.9 = \qty{-561.6}{kJ}$;
a \qty{1180.2}{K}: $-715.6 + 262.3 = -946.2 + 492.9 = \qty{-453.3}{kJ}$. Las
rectas se unen, como debe ser.
\textbf{8.} Pendiente $(-453.0 + 418.2)/200 = \qty{-0.174}{kJ/K}$, $a = -418.2 +
0.174 \times 1100 = \qty{-226.8}{kJ}$: $\Delta_r G^\circ = -226.8 - 0.174\,T$.
\textbf{9.} $\Delta_r S^\circ = \qty{+174}{J/K}$: un mol de gas da dos.
\textbf{10.} \ce{2ZnO + 2C -> 2Zn(g) + 2CO}, $\Delta_r G^\circ = (-226.8 - 0.174\,T)
- (-946.2 + 0.4177\,T) = 719.4 - 0.592\,T$ (\unit{kJ}).
\textbf{11.} Se anula en $T = 719.4/0.592 = \qty{1215}{K}$.
\textbf{12.} Por encima de su punto de ebullición el cinc se forma como vapor: la reacción produce
gas a partir de sólidos, su entropía es grande y positiva, y el vapor abandona la
carga, lo que separa el metal del mineral y del coque por destilación.
\textbf{13.} Cuatro especies, una reacción, una condición particular ($p_{\ce{Zn}} =
p_{\ce{CO}}$), tres fases: $v = 4 - 1 - 1 + 2 - 3 = 1$. Fijar la presión
fija la temperatura de equilibrio.
\textbf{14.} Un mol de vapor de cinc por mol de monóxido de carbono:
$p_{\ce{Zn}} = p_{\ce{CO}} = \qty{0.5}{bar}$.
\textbf{15.} Por mol de \ce{ZnO}, $\Delta_r G^\circ = \frac12(719.4 - 0.592 \times
1300) = \qty{-25.1}{kJ/mol}$, $K^\circ = \exp(25\,100/(8.314 \times 1300)) = 10$;
$Q = 0.25 < K^\circ$: la reducción avanza.
\textbf{16.} Justo por encima del cruce la reducción ya está favorecida (con
\qty{0.5}{bar} de cada gas, a partir de unos \qty{1170}{K}); más calor cuesta combustible,
desgasta las retortas de arcilla y no da más cinc, ya que la reacción llega
a ser completa de todos modos.
\textbf{17.} \ce{Zn(g) + CO2 -> ZnO + CO}.
\textbf{18.} La recta del óxido de cinc está por debajo de la recta \ce{CO}/\ce{CO2}
($-445$ frente a \qty{-357}{kJ} a \qty{1200}{K}): el vapor de cinc reduce el dióxido
de carbono y vuelve a convertirse en óxido. En la retorta el carbono caliente destruye todo
\ce{CO2} (Boudouard); en el condensador, más frío, nada lo hace, y el \ce{CO2}
---procedente de aire que entra por fugas, o de \ce{2CO -> C + CO2} al enfriarse--- convierte el cinc en un
polvo gris oxidado.
\textbf{19.} Un enfriamiento lento deja el vapor mucho tiempo en el intervalo en el que
se reoxida, y da un polvo fino en lugar de metal líquido; una condensación rápida
al estado líquido limita ambas cosas.
\textbf{20.} $10^6/65.38 = \qty{1.53e4}{mol}$ de cinc, y otros tantos de carbono:
$1.53 \times 10^4 \times 12.011 = \qty{184}{kg}$ (mucho más en la práctica, para calentar las
retortas).
\textbf{21.} $\Delta_r H^\circ = \frac12(-226.8 + 946.2) = \qty{+360}{kJ}$ por
mol de cinc; por tonelada, $1.53 \times 10^4 \times 360 = \qty{5.5e6}{kJ}$, es decir,
\qty{5.5}{GJ}, que se aportan quemando más carbón alrededor de las retortas.
\textbf{22.} El carbono reduce el óxido de cinc en condiciones estándar por encima de
$\boldsymbol{T \approx \qty{1.22e3}{K}}$.
\end{solution}
